\section{Related Work}
\label{sec:related}

\paragraph{Deep research systems.}
OpenAI Deep Research \citep{openai2025deepresearch} is an agent that plans searches, browses the web, and writes a cited report. It was trained with reinforcement learning on browsing tasks of the kind collected in BrowseComp \citep{wei2025browsecomp}, where each answer is a short fact that is difficult to find. Its training objective therefore rewards locating and citing information rather than structuring an argument, and its intermediate reasoning is held in a single context window with no explicit dependencies between sub-questions. STORM \citep{shao2024storm} writes Wikipedia-like articles. It first identifies perspectives on a topic, then simulates conversations in which writers holding those perspectives question an expert grounded in search results, and finally organizes the answers into an outline and an article. Because its questions arise from these conversations as the research proceeds, its process is bottom-up, and its output is deliberately balanced. Both systems were designed for summarization, and the comparison below treats them as the strongest available summary-oriented baselines rather than as argumentation systems.

\paragraph{Decomposition and structured reasoning.}
Breaking a problem into sub-problems before solving it is an established prompting strategy. Least-to-most prompting \citep{zhou2023leasttomost}, decomposed prompting \citep{khot2023decomposed}, and self-ask \citep{press2023selfask} each have a model generate sub-questions whose answers feed the original question. Tree of Thoughts \citep{yao2023tot} searches over intermediate reasoning steps, ReAct \citep{yao2023react} interleaves reasoning with calls to external tools, and Graph of Thoughts \citep{besta2024got} generalizes the tree to a graph in which intermediate results can merge. These methods target problems with verifiable answers, such as arithmetic, puzzles, and short-form question answering. They do not synthesize prose from retrieved evidence, and they do not revise their reasoning structure when a combined answer proves incomplete. \sys{} adopts the decomposition idea and extends it with explicit composition instructions, leaves grounded in retrieved evidence, and refinement of the graph when gaps appear.

\paragraph{Multi-agent frameworks.}
AutoGen \citep{wu2023autogen}, CAMEL \citep{li2023camel}, and natural-language societies of mind \citep{zhuge2023mindstorms} coordinate several language-model agents through conversation, assigning roles such as planner, researcher, and critic. In these frameworks the agents exchange messages, but the structure of the task itself is not represented, so the output of one agent does not formally constrain how another combines it. \sys{} makes that structure explicit in the DAG, so that the answer of each sub-question constrains the composition of its parent.

\paragraph{Argument mining and generation.}
Argument mining identifies claims, premises, and the relations between them in existing text \citep{lawrence2020argmining}. Argument generation produces new arguments. \citet{hua2018argument} generate counter-arguments with a sequence-to-sequence model conditioned on retrieved evidence, and Project Debater \citep{slonim2021debater} assembles debate speeches from a large mined corpus of claims and evidence. These systems draw on fixed corpora and generate text in a single pass. \sys{} instead retrieves evidence from the live web and constructs the structure of the argument before filling it in.

\paragraph{Citation and attribution.}
Text generated with citations can be checked for whether each citation supports the sentence it accompanies. \citet{rashkin2023ais} formalize this requirement as attribution to identified sources. \citet{liu2023verifiability} applied human judgments to four generative search engines and found that only about half of the generated sentences were fully supported by their citations, and that about three quarters of the citations supported their sentence. \citet{gao2023alce} built a benchmark that measures citation precision and recall automatically with an entailment model. These studies evaluate generation from retrieved passages in a single step. \Cref{sec:citations} examines a multi-step setting, in which citations must remain correct through several rounds of composition, and uses a lexical test that requires no additional model.

\paragraph{Implementation.}
The pipeline is written with DSPy \citep{khattab2024dspy}, which expresses each language-model call as a signature, consisting of inputs, outputs, and an instruction, and allows these calls to be composed into a program. The signatures used here are listed in \Cref{app:prompts}.
